\originalchapter{自由标量场与粒子解释}\label{ch:free-scalar}
\originalsection{模展开与正则量子化}
自由实标量的拉格朗日密度取 $\Lag_0=\frac12\partial_\mu\phi\partial^\mu\phi-\frac12m^2\phi^2$, 方程是 $(\Box+m^2)\phi=0$. 空间 Fourier 变换把它变成 $\ddot\phi_{\mathbf p}+E_{\mathbf p}^2\phi_{\mathbf p}=0$. 因而每个动量是频率为 $E_{\mathbf p}$ 的谐振子, 但实场的 $\mathbf p$ 与 $-\mathbf p$ 系数具有共轭关系, 不能再独立计数两次. 量子化后用厄米场算符写
\begin{equation}\label{eq:scalar-mode}
\phi(x)=\intp\frac1{\sqrt{2E_{\mathbf p}}}
\left(a_{\mathbf p}\e^{-\ii p\cdot x}+a^\dagger_{\mathbf p}\e^{\ii p\cdot x}\right),
\qquad p^0=E_{\mathbf p}.
\end{equation}
$a_{\mathbf p}$ 消灭一个动量 $\mathbf p$ 的激发, $a^\dagger_{\mathbf p}$ 产生一个. 两者在 Heisenberg 图景下是给定初始条件的时间无关模系数, 时间变化已放在平面波中. 本册采用 $\intp=\int\mathrm d^3p/(2\pi)^3$, $\delta^{(3)}$ 不含额外 $2\pi$ 因子.

量子化要求等时 $[\phi(t,\mathbf x),\pi(t,\mathbf y)]=\ii\delta^{(3)}(\mathbf x-\mathbf y)$, 同类场的等时交换子为零, 且 $\pi=\dot\phi$. 为什么模系数需带 $1/\sqrt{2E}$? 在 $t=0$ 写 $\widetilde\phi(\mathbf p)=\int\mathrm d^3x\,\e^{-\ii\mathbf p\cdot\mathbf x}\phi(0,\mathbf x)$, 就有
\begin{align*}
\widetilde\phi(\mathbf p)&=\frac{a_{\mathbf p}+a^\dagger_{-\mathbf p}}{\sqrt{2E_{\mathbf p}}},&
\widetilde\pi(\mathbf p)&=-\ii\sqrt{\frac{E_{\mathbf p}}2}(a_{\mathbf p}-a^\dagger_{-\mathbf p}),\\
a_{\mathbf p}&=\sqrt{\frac{E_{\mathbf p}}2}\widetilde\phi(\mathbf p)+
\frac\ii{\sqrt{2E_{\mathbf p}}}\widetilde\pi(\mathbf p).
\end{align*}
基本交换关系的 Fourier 变换为 $[\widetilde\phi(\mathbf p),\widetilde\pi(\mathbf q)]=\ii(2\pi)^3\delta^{(3)}(\mathbf p+\mathbf q)$. 代入逆式, 两个交叉项相加, 便得
\begin{equation}\label{eq:scalar-a-comm}
[a_{\mathbf p},a^\dagger_{\mathbf q}]=(2\pi)^3\delta^{(3)}(\mathbf p-\mathbf q),\qquad
[a_{\mathbf p},a_{\mathbf q}]=[a^\dagger_{\mathbf p},a^\dagger_{\mathbf q}]=0.
\end{equation}
反过来把此式代入场展开, 两个 $\phi\pi$ 交叉项各给 $\ii\delta^{(3)}/2$, 恢复正则交换子. 幺正时间演化保持算符乘积与交换子, 因而初始时刻施加的等时关系在以后各时刻继续成立.

\originalsection{能量、动量与真空}
将模展开代入 $H=\frac12\int\mathrm d^3x\,[\pi^2+(\nabla\phi)^2+m^2\phi^2]$, 空间积分把两动量通过 delta 关联. 两个湮灭算符的系数为 $-E_{\mathbf p}^2+\mathbf p^2+m^2=0$, 两个产生算符也消去. 混合项则留下
\[
H=\frac12\intp E_{\mathbf p}(a^\dagger_{\mathbf p}a_{\mathbf p}+a_{\mathbf p}a^\dagger_{\mathbf p})
=\intp E_{\mathbf p}a^\dagger_{\mathbf p}a_{\mathbf p}+E_0.
\]
在体积 $V$ 的周期边界盒中, $\mathbf p$ 离散, $E_0=\frac12\sum_{\mathbf p}E_{\mathbf p}$. 连续记号下 $E_0/V=\frac12\intp E_{\mathbf p}$ 紫外发散. 对平直时空内不涉及引力的自由粒子能量差, 可以减去这一定值; 正规序记号 $:H:$ 指把产生算符放在湮灭算符左侧, 在此自由情形正好实现该减法. 有边界时比较不同真空的能量差可能留下 Casimir 效应, 所以不能把一切真空效应都因正规序而宣布消失.

动量同样得到 $\mathbf P=\intp\mathbf p\,a^\dagger_{\mathbf p}a_{\mathbf p}$; 零点项因动量反演对称而为零. 对任意模有 $[H,a^\dagger_{\mathbf p}]=E_{\mathbf p}a^\dagger_{\mathbf p}$, $[\mathbf P,a^\dagger_{\mathbf p}]=\mathbf p\,a^\dagger_{\mathbf p}$. 这两式给出了把场激发解释为粒子的直接依据, 而不是先假设有粒子再附上名称.

\begin{definition}
真空 $\ket0$ 是满足 $a_{\mathbf p}\ket0=0$ 的归一态. 对有限体积的离散模, $\ket{n_{\mathbf p}}=(a^\dagger_{\mathbf p})^n\ket0/\sqrt{n!}$ 为数态. 全部有限粒子态及其完备化构成玻色 Fock 空间 $\mathcal F=\bigoplus_{N\ge0}\operatorname{Sym}\mathcal H_1^{\otimes N}$.
\end{definition}
其中 $\mathcal H_1$ 是一粒子态空间, $\operatorname{Sym}$ 表示交换粒子标签后状态不变. 两个产生算符可交换, 故同一模的占据数可以是任意非负整数. 连续动量本征态不是普通可归一向量, 必须按分布处理. 取相对论归一 $\ket{\mathbf p}=\sqrt{2E_{\mathbf p}}a^\dagger_{\mathbf p}\ket0$, 则
\[
\braket{\mathbf q}{\mathbf p}=2E_{\mathbf p}(2\pi)^3\delta^{(3)}(\mathbf p-\mathbf q),\qquad
1_{\mathcal H_1}=\intp\frac{\ket{\mathbf p}\bra{\mathbf p}}{2E_{\mathbf p}}.
\]
质量壳测度的不变性解释为什么这套归一在散射论中更方便. 波包 $\ket f=\intp f(\mathbf p)\ket{\mathbf p}/\sqrt{2E_{\mathbf p}}$ 满足 $\langle f|f\rangle=\intp|f(\mathbf p)|^2$, 从而成为真正可归一态.

\originalsection{复场与反粒子}
复自由场使用两个独立的模族:
\begin{equation}\label{eq:complex-mode}
\phi(x)=\intp\frac1{\sqrt{2E_{\mathbf p}}}
\left(a_{\mathbf p}\e^{-\ii p\cdot x}+b^\dagger_{\mathbf p}\e^{\ii p\cdot x}\right).
\end{equation}
其共轭场为 $\phi^\dagger(x)=\intp(a^\dagger_{\mathbf p}\e^{\ii p\cdot x}+b_{\mathbf p}\e^{-\ii p\cdot x})/\sqrt{2E_{\mathbf p}}$. 正则关系 $[\phi(\mathbf x),\dot\phi^\dagger(\mathbf y)]=\ii\delta^{(3)}(\mathbf x-\mathbf y)$ 及共轭关系要求 $[a,a^\dagger]$ 和 $[b,b^\dagger]$ 均为式\eqref{eq:scalar-a-comm}, 交叉交换子为零. 因而正规序后的能量和电荷分别为
\begin{align*}
:H:&=\intp E_{\mathbf p}(a^\dagger_{\mathbf p}a_{\mathbf p}+b^\dagger_{\mathbf p}b_{\mathbf p}),\\
:Q:&=\intp(a^\dagger_{\mathbf p}a_{\mathbf p}-b^\dagger_{\mathbf p}b_{\mathbf p}).
\end{align*}
电荷来自上一章约定的 $j^0=\ii(\phi^\dagger\dot\phi-\dot\phi^\dagger\phi)$. 具体代入时 $a^\dagger a$ 的系数是 $+1$, $bb^\dagger$ 是 $-1$; 后者正规序后给 $-b^\dagger b$ 并去除真空常数. 能量却对两族都为正. 所以负频部分创建反粒子, 并不是创建负能粒子.

由 $[Q,a]=-a$, $[Q,b^\dagger]=-b^\dagger$ 得 $[Q,\phi]=-\phi$. 对称性用 $U(\alpha)=\e^{\ii\alpha Q}$ 实现, 有 $U\phi U^{-1}=\e^{-\ii\alpha}\phi$, 与上一章的相位约定相符. $a^\dagger\ket0$ 和 $b^\dagger\ket0$ 分别带 $+1,-1$ 电荷, 但若把相位变换约定整体反向, 粒子电荷名称也随之反向. 相对电荷和电荷守恒不受这种命名影响.

\originalsection{场产生的态与位置解释}
在实场情形,
\[
\phi(x)\ket0=\intp\frac{\e^{\ii p\cdot x}}{\sqrt{2E_{\mathbf p}}}a^\dagger_{\mathbf p}\ket0.
\]
它是所有动量的一粒子叠加, 而 $\bra0\phi(x)\ket{\mathbf p}=\e^{-\ii p\cdot x}$ 正好是质量壳上的平面波. 这个矩阵元把场与熟悉的波函数联系起来, 但不能把 $\phi(x)\ket0$ 直接当作位置本征态: 它在一点处没有有限范数, 不同点的该类态也不是正交 delta 归一. 局域场的意义通过局域耦合和因果交换关系表达, 不等于非相对论位置算符的简单延续.

取外源 $J(x)$ 与场相互作用 $S_J=\int\mathrm d^4x\,J\phi$. 它描述仪器在局部位置对场作激发. 一阶微扰下, 在 $a^\dagger_{\mathbf p}\ket0$ 这套非协变态上的振幅为 $\ii\widetilde J(p)/\sqrt{2E_{\mathbf p}}$; 投影到本章定义的相对论归一态 $\ket{\mathbf p}$ 时则为 $\ii\widetilde J(p)$, 其中 $\widetilde J(p)=\int\mathrm d^4x\,\e^{\ii p\cdot x}J(x)$. 两种写法因末态归一不同而差一个 $\sqrt{2E}$, 计算概率必须同时使用对应完备测度. 仪器的时空形状决定被激发的动量波包; 单个局域点不提供独立可归一粒子态. 下一章利用两次场作用的真空矩阵元描述此类激发之间的相关.
